% GENERATED FILE. Edit canonical lessons, then run npm run generate:books.
% canonical-source-hash: fnv1a64:441a36aca0022695
% canonical-lessons: PA-C301-hal, PA-C301-balad, PA-C301-jula, PA-C301-vahuna, PA-C301-jotna, PA-R301-review-295-301

\chapter{Ploughing the Field}
\label{ch:pa-a2-ploughing-the-field}
\hlchaptermodality{301}

\begin{chapteropening}
\textbf{By the end of this chapter:} \emph{I can talk about ploughing with oxen.}

Complete the last lesson of chapter 301: I can talk about ploughing with oxen.
\end{chapteropening}

\section[\texorpdfstring{\pa{ਹਲ} (hal)}{ਹਲ (hal)}]{\pa{ਹਲ} (hal) — a plough}

\label{lesson:PA-C301-hal}



\begin{quote}
Before the new one: say the Punjabi for raw cotton, then the Punjabi for sugar cane.
\end{quote}

\subsection*{You'll want to know: \pa{ਹਲ}}
\textbf{\pa{ਹਲ}} — \emph{hal} — \textquotedblleft{}a plough\textquotedblright{}.

\begin{quote}
\pa{ਸਾਡਾ} \pa{ਹਲ} \pa{ਪੁਰਾਣਾ} \pa{ਹੈ।}
\end{quote}

\begin{quote}
\emph{sāḍā hal purāṇā hai.}
\end{quote}

\begin{quote}
Our plough is old.
\end{quote}

A masculine noun.

\begin{cousinweb}
\textbf{\pa{ਹਲ}} (\emph{hal}) is inherited from Sanskrit \emph{hala}, \textquotedblleft{}plough\textquotedblright{}; CDIAL lists Punjabi \emph{hal} under it (CDIAL 14000).
\end{cousinweb}

\begin{grammarlens}[title={What you've built}]
One more word for this chapter.
\end{grammarlens}

\subsection*{Your turn}
\begin{itemize}
\raggedright
  \item \emph{Say it:} \emph{hal}
  \item \emph{Say it:} \emph{hal}, once more
  \item \emph{Say it:} \emph{sāḍā hal purāṇā hai.} — Our plough is old.
  \item \emph{Recall:} say the Punjabi for raw cotton, then the Punjabi for sugar cane, then say \emph{hal} again
\end{itemize}

\begin{tcolorbox}[breakable,colback=teal!4,colframe=teal!35!black,title={Before you move on}]
What does \pa{ਹਲ} mean? (\textquotedblleft{}A plough\textquotedblright{}.) Say it once more.
\end{tcolorbox}
\section[\texorpdfstring{\pa{ਬਲਦ} (balad)}{ਬਲਦ (balad)}]{\pa{ਬਲਦ} (balad) — an ox; a bullock}

\label{lesson:PA-C301-balad}



\begin{quote}
Before the new one: say the Punjabi for sugar cane, then the Punjabi for a plough.
\end{quote}

\subsection*{You'll want to know: \pa{ਬਲਦ}}
\textbf{\pa{ਬਲਦ}} — \emph{balad} — \textquotedblleft{}an ox; a bullock\textquotedblright{}.

\begin{quote}
\pa{ਸਾਡੇ} \pa{ਕੋਲ} \pa{ਦੋ} \pa{ਬਲਦ} \pa{ਹਨ।}
\end{quote}

\begin{quote}
\emph{sāḍe kol do balad han.}
\end{quote}

\begin{quote}
We have two oxen.
\end{quote}

A masculine noun.

\begin{cousinweb}
\textbf{\pa{ਬਲਦ}} (\emph{balad}) is inherited from Sanskrit \emph{balivarda}, \textquotedblleft{}ox, bull\textquotedblright{}; CDIAL lists Punjabi \emph{bald}, \textquotedblleft{}ox\textquotedblright{}, under it (CDIAL 9176).
\end{cousinweb}

\begin{grammarlens}[title={What you've built}]
One more word for this chapter.
\end{grammarlens}

\subsection*{Your turn}
\begin{itemize}
\raggedright
  \item \emph{Say it:} \emph{balad}
  \item \emph{Say it:} \emph{balad}, once more
  \item \emph{Say it:} \emph{sāḍe kol do balad han.} — We have two oxen.
  \item \emph{Recall:} say the Punjabi for sugar cane, then the Punjabi for a plough, then say \emph{balad} again
  \item \emph{Write it:} copy \textbf{\pa{ਬਲਦ}} once, checking the whole word
\end{itemize}

\begin{tcolorbox}[breakable,colback=teal!4,colframe=teal!35!black,title={Before you move on}]
What does \pa{ਬਲਦ} mean? (\textquotedblleft{}An ox; a bullock\textquotedblright{}.) Say it once more.
\end{tcolorbox}
\section[\texorpdfstring{\pa{ਜੂਲਾ} (jūlā)}{ਜੂਲਾ (jūlā)}]{\pa{ਜੂਲਾ} (jūlā) — a yoke}

\label{lesson:PA-C301-jula}



\begin{quote}
Before the new one: say the Punjabi for a plough, then the Punjabi for an ox.
\end{quote}

\subsection*{You'll want to know: \pa{ਜੂਲਾ}}
\textbf{\pa{ਜੂਲਾ}} — \emph{jūlā} — \textquotedblleft{}a yoke\textquotedblright{}.

\begin{quote}
\pa{ਜੂਲਾ} \pa{ਬਲਦਾਂ} \pa{ਦੀ} \pa{ਧੌਣ} \pa{ਉੱਤੇ} \pa{ਹੁੰਦਾ} \pa{ਹੈ।}
\end{quote}

\begin{quote}
\emph{jūlā baldāṁ dī dhauṇ utte hundā hai.}
\end{quote}

\begin{quote}
The yoke sits on the oxen's necks.
\end{quote}

A masculine noun.

\begin{cousinweb}
\textbf{\pa{ਜੂਲਾ}} (\emph{jūlā}) is inherited from a reconstructed \emph{*yugahala}, \textquotedblleft{}plough with yoke\textquotedblright{}; CDIAL lists Punjabi \emph{jūlā}, \textquotedblleft{}crossbar of yoke\textquotedblright{}, under it (CDIAL 10491).
\end{cousinweb}

\begin{grammarlens}[title={What you've built}]
One more word for this chapter.
\end{grammarlens}

\subsection*{Your turn}
\begin{itemize}
\raggedright
  \item \emph{Say it:} \emph{jūlā}
  \item \emph{Say it:} \emph{jūlā}, once more
  \item \emph{Say it:} \emph{jūlā baldāṁ dī dhauṇ utte hundā hai.} — The yoke sits on the oxen's necks.
  \item \emph{Recall:} say the Punjabi for a plough, then the Punjabi for an ox, then say \emph{jūlā} again
  \item \emph{Write it:} copy \textbf{\pa{ਜੂਲਾ}} once, checking the whole word
\end{itemize}

\begin{tcolorbox}[breakable,colback=teal!4,colframe=teal!35!black,title={Before you move on}]
What does \pa{ਜੂਲਾ} mean? (\textquotedblleft{}A yoke\textquotedblright{}.) Say it once more.
\end{tcolorbox}
\section[\texorpdfstring{\pa{ਵਾਹੁਣਾ} (vāhuṇā)}{ਵਾਹੁਣਾ (vāhuṇā)}]{\pa{ਵਾਹੁਣਾ} (vāhuṇā) — to plough}

\label{lesson:PA-C301-vahuna}



\begin{quote}
Before the new one: say the Punjabi for an ox, then the Punjabi for a yoke.
\end{quote}

\subsection*{You'll want to know: \pa{ਵਾਹੁਣਾ}}
\textbf{\pa{ਵਾਹੁਣਾ}} — \emph{vāhuṇā} — \textquotedblleft{}to plough\textquotedblright{}.

\begin{quote}
\pa{ਅਸੀਂ} \pa{ਕੱਲ੍ਹ} \pa{ਖੇਤ} \pa{ਵਾਹਾਂਗੇ।}
\end{quote}

\begin{quote}
\emph{asīṁ kall khet vāhāṅge.}
\end{quote}

\begin{quote}
We will plough the field tomorrow.
\end{quote}

Past (I …): \pa{ਮੈਂ} \pa{ਖੇਤ} \pa{ਵਾਹਿਆ}, \emph{maiṁ khet vāhiā}.

\begin{cousinweb}
\textbf{\pa{ਵਾਹੁਣਾ}} (\emph{vāhuṇā}) is inherited from Sanskrit \emph{vāhayati}, \textquotedblleft{}drives (a chariot)\textquotedblright{}; CDIAL lists Punjabi \emph{vāhuṇā}, \textquotedblleft{}to drive (a beast), plough\textquotedblright{}, under it (CDIAL 11612).
\end{cousinweb}

\begin{grammarlens}[title={What you've built}]
One more word for this chapter.
\end{grammarlens}

\subsection*{Your turn}
\begin{itemize}
\raggedright
  \item \emph{Say it:} \emph{vāhuṇā}
  \item \emph{Say it:} \emph{vāhuṇā}, once more
  \item \emph{Say it:} \emph{asīṁ kall khet vāhāṅge.} — We will plough the field tomorrow.
  \item \emph{Recall:} say the Punjabi for an ox, then the Punjabi for a yoke, then say \emph{vāhuṇā} again
  \item \emph{Write it:} copy \textbf{\pa{ਵਾਹੁਣਾ}} once, checking the whole word
\end{itemize}

\begin{tcolorbox}[breakable,colback=teal!4,colframe=teal!35!black,title={Before you move on}]
What does \pa{ਵਾਹੁਣਾ} mean? (\textquotedblleft{}To plough\textquotedblright{}.) Say it once more.
\end{tcolorbox}
\section[\texorpdfstring{\pa{ਜੋਤਣਾ} (jotṇā)}{ਜੋਤਣਾ (jotṇā)}]{\pa{ਜੋਤਣਾ} (jotṇā) — to yoke; to harness}

\label{lesson:PA-C301-jotna}



\begin{quote}
Before the new one: say the Punjabi for a yoke, then the Punjabi for to plough.
\end{quote}

\subsection*{You'll want to know: \pa{ਜੋਤਣਾ}}
\textbf{\pa{ਜੋਤਣਾ}} — \emph{jotṇā} — \textquotedblleft{}to yoke; to harness\textquotedblright{}.

\begin{quote}
\pa{ਮੈਂ} \pa{ਸਵੇਰੇ} \pa{ਬਲਦ} \pa{ਜੋਤਦਾ} \pa{ਹਾਂ।}
\end{quote}

\begin{quote}
\emph{maiṁ savere balad jotdā hāṁ.}
\end{quote}

\begin{quote}
I yoke the oxen in the morning.
\end{quote}

Past (I …): \pa{ਮੈਂ} \pa{ਬਲਦ} \pa{ਜੋਤੇ}, \emph{maiṁ balad jote}.

\begin{cousinweb}
\textbf{\pa{ਜੋਤਣਾ}} (\emph{jotṇā}) is inherited from Sanskrit \emph{yōktrayati}, \textquotedblleft{}fastens, ties\textquotedblright{}; CDIAL lists Punjabi \emph{jotṇā}, \textquotedblleft{}to yoke, plough\textquotedblright{}, under it (CDIAL 10524).
\end{cousinweb}

\begin{grammarlens}[title={What you've built}]
One more word for this chapter.
\end{grammarlens}

\subsection*{Your turn}
\begin{itemize}
\raggedright
  \item \emph{Say it:} \emph{jotṇā}
  \item \emph{Say it:} \emph{jotṇā}, once more
  \item \emph{Say it:} \emph{maiṁ savere balad jotdā hāṁ.} — I yoke the oxen in the morning.
  \item \emph{Recall:} say the Punjabi for a yoke, then the Punjabi for to plough, then say \emph{jotṇā} again
  \item \emph{Write it:} copy \textbf{\pa{ਜੋਤਣਾ}} once, checking the whole word
\end{itemize}

\begin{tcolorbox}[breakable,colback=teal!4,colframe=teal!35!black,title={Before you move on}]
What does \pa{ਜੋਤਣਾ} mean? (\textquotedblleft{}To yoke; to harness\textquotedblright{}.) Say it once more.
\end{tcolorbox}
\section[(dialogue)]{Looking back: from Cracks and Hollows to Ploughing the Field}

\label{lesson:PA-R301-review-295-301}



\begin{quote}
Say the Punjabi for a crack and to split open.
\end{quote}

\subsection*{Your turn}
Say these aloud:

\begin{itemize}
\raggedright
  \item a crack, to split open, hollow, a ravine, to rot
  \item to slide, to stick, thick, slick, to melt
  \item smooth, fine-grained, tender, blunt, to prick
  \item a thorn, a bush, an acacia tree, dense, to wilt
  \item a riverbank, a field channel, a wave, brackish, beautiful
  \item a crop, the spring harvest, the autumn harvest, raw cotton, sugar cane
  \item a plough, an ox, a yoke, to plough, to yoke
\end{itemize}

\begin{tcolorbox}[breakable,colback=teal!4,colframe=teal!35!black,title={Before you move on}]
A crack? (\textbf{\pa{ਤਰੇੜ}}, \emph{tareṛ}.) To split open? (\textbf{\pa{ਫਟਣਾ}}, \emph{phaṭṇā}.) Hollow? (\textbf{\pa{ਖੋਖਲਾ}}, \emph{khokhlā}.) A ravine? (\textbf{\pa{ਖੱਡ}}, \emph{khaḍḍ}.) To rot? (\textbf{\pa{ਗਲਣਾ}}, \emph{galṇā}.) To slide? (\textbf{\pa{ਖਿਸਕਣਾ}}, \emph{khisakṇā}.) To stick? (\textbf{\pa{ਚਿਪਕਣਾ}}, \emph{chipakṇā}.) Thick? (\textbf{\pa{ਗਾੜ੍ਹਾ}}, \emph{gāṛhā}.) Slick? (\textbf{\pa{ਚਿਕਨਾ}}, \emph{chiknā}.) To melt? (\textbf{\pa{ਪਿਘਲਣਾ}}, \emph{pighalṇā}.) Smooth? (\textbf{\pa{ਮੁਲਾਇਮ}}, \emph{mulāim}.) Fine-grained? (\textbf{\pa{ਬਰੀਕ}}, \emph{barīk}.) Tender? (\textbf{\pa{ਕੂਲਾ}}, \emph{kūlā}.) Blunt? (\textbf{\pa{ਖੁੰਢਾ}}, \emph{khuṇḍhā}.) To prick? (\textbf{\pa{ਚੁਭਣਾ}}, \emph{chubhṇā}.) A thorn? (\textbf{\pa{ਕੰਡਾ}}, \emph{kaṇḍā}.) A bush? (\textbf{\pa{ਬੂਟਾ}}, \emph{būṭā}.) An acacia tree? (\textbf{\pa{ਕਿੱਕਰ}}, \emph{kikkar}.) Dense? (\textbf{\pa{ਸੰਘਣਾ}}, \emph{saṅghṇā}.) To wilt? (\textbf{\pa{ਕੁਮਲਾਉਣਾ}}, \emph{kumlāuṇā}.) A riverbank? (\textbf{\pa{ਕੰਢਾ}}, \emph{kaṇḍhā}.) A field channel? (\textbf{\pa{ਖਾਲ}}, \emph{khāl}.) A wave? (\textbf{\pa{ਲਹਿਰ}}, \emph{lahir}.) Brackish? (\textbf{\pa{ਖਾਰਾ}}, \emph{khārā}.) Beautiful? (\textbf{\pa{ਸੋਹਣਾ}}, \emph{sohṇā}.) A crop? (\textbf{\pa{ਫ਼ਸਲ}}, \emph{fasal}.) The spring harvest? (\textbf{\pa{ਹਾੜ੍ਹੀ}}, \emph{hāṛhī}.) The autumn harvest? (\textbf{\pa{ਸਾਉਣੀ}}, \emph{sāuṇī}.) Raw cotton? (\textbf{\pa{ਕਪਾਹ}}, \emph{kapāh}.) Sugar cane? (\textbf{\pa{ਗੰਨਾ}}, \emph{gannā}.) A plough? (\textbf{\pa{ਹਲ}}, \emph{hal}.) An ox? (\textbf{\pa{ਬਲਦ}}, \emph{balad}.) A yoke? (\textbf{\pa{ਜੂਲਾ}}, \emph{jūlā}.) To plough? (\textbf{\pa{ਵਾਹੁਣਾ}}, \emph{vāhuṇā}.) To yoke? (\textbf{\pa{ਜੋਤਣਾ}}, \emph{jotṇā}.)
\end{tcolorbox}
